% !TEX root = ../article.tex
% =============================================================================
% 3.4 Datos y preparación
% =============================================================================
% Guía: origen, tamaño y características de los datos; criterios de
% selección, limpieza, etiquetado y preprocesamiento; separación de
% entrenamiento, validación y prueba (con modelos preentrenados, un conjunto
% de evaluación separado de los ejemplos usados para ajustar la solución);
% limitaciones de calidad o cobertura y tratamiento de información sensible.
%
% Primer avance: datos disponibles o fuente prevista y cómo se accederá a
% ellos.
%
% Datos de cada conjunto verificados el 8 de noviembre de 2026 en Zenodo
% (API de registros), en la página de la colección en TCIA, en el artículo de
% cada conjunto y en el repositorio del benchmark de COVID-19.
% =============================================================================

\section{Datos y preparación}\label{sec:data}

\subsection{Origen y características}\label{sec:data-source}
Thoraxis usa tres conjuntos públicos de TC de tórax con máscaras de referencia
de ambos pulmones, todos anonimizados por sus autores
(Tabla~\ref{tab:datasets}).

\begin{itemize}
    \item \textbf{TotalSegmentator}, versión~3: 1939 TC tomadas al azar de la
          rutina clínica y reunidas por el Hospital Universitario de Basilea,
          con una amplia variedad de
          patologías, equipos, secuencias e instituciones, y 117 estructuras
          segmentadas \parencite{wasserthal2023dataset}. La versión~3 agregó
          291 estudios pediátricos. Los pulmones están segmentados por
          lóbulos, con la tráquea como clase aparte \parencite{wasserthal2023}.
    \item \textbf{COVID-19 CT Lung and Infection Segmentation}: 20 TC de
          pacientes con COVID-19 de la Coronacases Initiative y de
          Radiopaedia, con el pulmón izquierdo, el derecho y la infección
          delimitados corte por corte \parencite{ma2020}. Anotadores con uno a
          cinco años de experiencia trazaron las máscaras, dos radiólogos las
          refinaron y un radiólogo con más de diez años de experiencia en
          tórax las verificó \parencite{ma2021}.
    \item \textbf{LCTSC} (\textit{Lung CT Segmentation Challenge} 2017): 60
          TC de planificación de radioterapia de pacientes con cáncer de
          pulmón, 20 de cada una de tres instituciones, con los contornos
          clínicos de pulmón izquierdo, pulmón derecho, corazón, esófago y
          médula espinal, revisados según la guía RTOG~1106
          \parencite{yang2017, yang2018}.
\end{itemize}

\begin{table}[htbp]
    \centering
    \caption{Conjuntos de datos de Thoraxis.}\label{tab:datasets}
    \small
    \begin{tabularx}{\textwidth}{>{\raggedright\arraybackslash}p{3cm}
                                 >{\raggedright\arraybackslash}p{1.4cm}
                                 >{\raggedright\arraybackslash}X
                                 >{\raggedright\arraybackslash}X
                                 >{\raggedright\arraybackslash}p{2.7cm}}
        \toprule
        Conjunto & Estudios & Formato, tamaño y acceso & Licencia & Uso en
            Thoraxis \\
        \midrule
        TotalSegmentator v3 & 1939 & 37,4~GB en Zenodo, descarga directa. &
            CC BY 4.0. & Entrenamiento, validación y prueba interna. \\
        COVID-19 CT Lung and Infection & 20 & NIfTI; 1,1~GB en Zenodo,
            descarga directa. & CC BY 4.0 en Zenodo; CC BY-NC-SA en el artículo y el
            repositorio. & Entrenamiento, validación y prueba interna. \\
        LCTSC & 60 & DICOM con RTSTRUCT; 5,15~GB en TCIA, con TCIA Data
            Retriever. & CC BY 3.0. & Solo validación externa. \\
        \bottomrule
    \end{tabularx}
    \articleSource{Elaboración propia a partir de
    \textcite{wasserthal2023dataset}, \textcite{ma2020} y
    \textcite{yang2017}.}
\end{table}

Los datos no se guardan en los repositorios. El repositorio
\texttt{thoraxis-model} contiene los scripts de descarga desde las fuentes
oficiales y versiona los datos procesados con DVC, de modo que cada integrante
obtiene los mismos archivos. Para explorar los datos antes de la descarga
completa, TotalSegmentator ofrece un subconjunto de 102 estudios
\parencite{wasserthal2023dataset}.

\subsection{Selección, limpieza y preprocesamiento}\label{sec:preprocessing}
De TotalSegmentator se seleccionan los estudios que contienen ambos pulmones
completos, es decir, con los cinco lóbulos presentes y sin cortes truncados en
los vértices ni en las bases. Se excluyen los estudios con archivos dañados o
con máscaras incompletas. El preprocesamiento sigue estos pasos:

\begin{enumerate}
    \item \textbf{Conversión:} todos los estudios se llevan a NIfTI; los de
          LCTSC se convierten desde DICOM y RTSTRUCT.
    \item \textbf{Etiquetas comunes:} cada estudio queda con dos clases,
          pulmón izquierdo y pulmón derecho. En TotalSegmentator se unen los
          lóbulos de cada lado; en LCTSC se usan \texttt{Lung\_L} y
          \texttt{Lung\_R}; en el conjunto de COVID-19 se descarta la máscara
          de infección, que queda incluida en el pulmón.
    \item \textbf{Orientación y remuestreo:} todos los volúmenes se llevan a
          la misma orientación y a un tamaño de vóxel común.
    \item \textbf{Intensidades:} se aplica la ventana pulmonar de $-1250$ a
          $250$ unidades Hounsfield y se normaliza a $[0, 1]$. Diez casos del
          conjunto de COVID-19, los de Radiopaedia, ya vienen con esa ventana
          y reescalados a $[0, 255]$, y sus autores recomiendan tratar los
          demás igual \parencite{ma2020}; aplicar la misma ventana a todos los
          conjuntos los vuelve comparables.
\end{enumerate}

\pendiente{tamaño de vóxel común, número de estudios de TotalSegmentator que
cumplen los criterios y decisión sobre los estudios pediátricos}

\subsection{Partición de los datos}\label{sec:data-split}
La partición se hace por paciente, nunca por corte, para que ningún paciente
aparezca en dos conjuntos. Los estudios de TotalSegmentator y de COVID-19 se
dividen en entrenamiento, validación y prueba interna, estratificando por
fuente para que los pulmones con COVID-19 estén en los tres conjuntos. LCTSC
queda completo como validación externa: sus 60 pacientes provienen de otras
fuentes y no se usan para entrenar ni para elegir hiperparámetros.

TotalSegmentator, el modelo publicado que sirve de referencia, se entrenó con
los estudios de su propio conjunto \parencite{wasserthal2023}. Por eso solo se
compara con Thoraxis en LCTSC y en los casos de COVID-19, que no vio durante
su entrenamiento.

\pendiente{proporciones de la partición y número de pacientes por conjunto}

\subsection{Limitaciones y datos sensibles}\label{sec:data-limitations}
Los conjuntos difieren en cómo delimitan el pulmón. En LCTSC, los hilios, los
vasos de más de 5~mm y los bronquios principales quedan fuera del pulmón, el
tumor se excluye en la mayoría de los casos sin garantía sobre la extensión
excluida, y el pulmón colapsado puede excluirse en algunos estudios
\parencite{yang2017}. Estas diferencias pueden bajar el Dice en la validación
externa aunque la segmentación sea anatómicamente razonable, y se analizarán
en la evaluación. Como referencia de lo alcanzable, en el reto de 2017 los
mejores métodos obtuvieron un Dice promedio de 0,95 a 0,98 para los pulmones
\parencite{yang2018}.

La cobertura también es limitada: el conjunto de COVID-19 tiene solo 20
casos, LCTSC representa únicamente a pacientes con cáncer de pulmón de tres
instituciones, y ninguno de los conjuntos incluye pacientes de la región. La
licencia del
conjunto de COVID-19 es ambigua, pues el registro de Zenodo indica CC BY 4.0
y el artículo y el repositorio del benchmark, CC BY-NC-SA; el proyecto se
ajusta a la más restrictiva, que permite su uso académico no comercial.

Todos los datos son públicos y anonimizados. Ningún estudio, máscara ni dato de
pacientes se sube a los repositorios, cuyo \texttt{.gitignore} debe bloquear
los formatos médicos (DICOM, NIfTI, NRRD y MHA); se citan todas las fuentes
según sus licencias y no se intenta identificar a ningún paciente.
